\makeatletter
\def\input@path{{./}{../}{../../}{../../../}{../../../../}{preamble/}{../preamble/}{../../preamble/}{../../../preamble/}}
\makeatother

\input{preamble_exams}
\input{graphs/G11_Histograms.tex}

\graphicspath{{./}{../}{../../}{../../../}{../../../../}{preamble/}{../preamble/}{../../preamble/}{../../../preamble/}}

% Learning Objective Column width
\setYColumnWidth{0.30}




\begin{document}
% HEADER BOX (3 COLUMNS)
\ExamHeader{11th}{C5, C6 and C7}{Continuous distributions and sample data}{1st}
\vspace{0.15cm}
% OBJECTIVE + CRITERIA
\LearningObjective
{Use continuous probability models and sample data to describe typical values,
variability, and interval probabilities, and interpret differences between
sample estimates and theoretical values.}
{C5: Calculates and interprets the mean, variance and standard deviation of continuous distributions using standard formulas and functions.\par
C6: Estimates interval probabilities using relative frequencies from generated samples.
 and compares them with theoretical probabilities.\par
C7: Calculates sample means, variances, and standard deviations for different sample sizes and data representations.}
% STUDENT LINE
\vspace{0.1cm}
% EXAM BODY
\begin{multicols}{2}
\SubsectionBox{Criteria assessment}\vspace{-0.25cm}
This exam assesses C5, C6 and C7 separately, with 24 points per criterion
(72 points in total). Each task shows its point allocation. Partial credit
may be awarded when the correct procedure is clearly shown but arithmetic
errors are made. Answers without the required procedure will not receive credit.
Failure to write your name on the task sheet or submitting it after the deadline
will result in a one-point deduction from each criterion.
For this exam, a minimum of 15 out of 24 points is required to pass each criterion.
Use unrounded intermediate values; round final statistics to three decimal
places and percentage errors to two decimal places where necessary.

\SubsectionBox{1. Two packing stations \score{C5-\sml{2}}}
A gift company uses two packing stations. At station A, an automated machine
finishes a gift box in anywhere from 4 to 16 minutes. Within these limits,
any two time windows of the same length are equally likely to contain the
finishing time. Let $A$ be its packing time in minutes.

At station B, staff pack custom gift boxes, also taking between 4 and 16
minutes. To describe many orders, the manager draws a smooth outline of how
often different packing times occur. The outline starts at zero at 4 minutes,
rises along a straight edge to its single highest point at 7 minutes, then
falls along a straight edge to zero at 16 minutes. No packing times occur
outside these limits. Let $B$ be this packing time in minutes.
\begin{enumerate}[label=\textbf{(\alph*)}]
\item Identify the distribution for A from the clues and calculate its mean, variance, and standard deviation using the corresponding formulas. \score{1}
\item Identify the distribution for B from the clues and calculate its mean, variance, and standard deviation using the corresponding formulas. \score{1}
\end{enumerate}

\SubsectionBox{2. Two custom-order workshops \score{C5-\sml{2}}}
Two workshops make custom wooden signs. Each order takes between 2 and
11 hours to complete. The owners use smooth outlines to describe how often
different completion times occur over many orders; times outside these limits
never occur.

At workshop I, longer jobs are more common: the outline is a single straight
edge rising from zero at 2 hours to its highest point at 11 hours.
At workshop D, shorter jobs are more common: the outline is a single straight
edge falling from its highest point at 2 hours to zero at 11 hours.
Let $I$ and $D$ be the respective completion times in hours.
\begin{enumerate}[label=\textbf{(\alph*)}]
\item Identify each distribution from the clues and calculate its mean, variance, and standard deviation using the corresponding formulas. \score{2}
\end{enumerate}

\newpage

\SubsectionBox{3. Arrivals and daily revenue \score{C5-\sml{2}}}
A bicycle repair shop receives online service requests throughout its opening
hours, averaging 8 requests per hour. Requests arrive independently, one at
a time, and the arrival pace stays the same throughout the day. In any very
short time window, the chance of a request is proportional to the window's
length, while two or more requests in that window are negligible. How long
the shop has already waited does not affect the chance of a request during
the next minute. Let $T$ be the time between consecutive requests in minutes.

On comparable working days, the shop's total revenue varies because of many
small, independent influences, with none dominating. A smooth outline of
daily revenues has a symmetric bell shape, highest at \$2400 and tapering
on both sides. The two places where its sides change from becoming steeper
to becoming flatter are \$2250 and \$2550. Treat this bell-shaped pattern
as exact for the calculations. Let $R$ be daily revenue measured in hundreds
of dollars.
\begin{enumerate}[label=\textbf{(\alph*)}]
\item Identify the distribution for T from the clues and use its formulas to find the mean, variance, and standard deviation in minutes. \score{1}
\item Identify the revenue distribution from the clues. Use its formulas to find the mean, variance, and standard deviation in dollars. Interpret the mean. \score{1}
\end{enumerate}

\SubsectionBox{4. Water consumption, raw data \score{C6-\sml{1}}}
A simulated sample of daily household water consumption (hundreds of liters) is
\[
\begin{array}{cccccccccc}
2.1&2.5&2.9&3.0&3.2&3.6&3.9&4.0&4.3&4.7\\
2.4&2.8&3.1&3.4&3.7&4.0&4.1&4.5&4.8&5.2
\end{array}
\]
The theoretical probability is $P(3\le X\le4)=0.50$.
\begin{enumerate}[label=\textbf{(\alph*)}]
\item Find the empirical probability of the stated interval from the given sample. \score{0.5}
\item Calculate the absolute and percentage errors relative to the given theoretical probability. \score{0.5}
\end{enumerate}

\SubsectionBox{5. Service times, frequency data \score{C6-\sml{1}}}
Recorded service times are summarized below.
\[
\begin{array}{c|rrrrrr}
 x_i\text{ (min)}&3&6&9&12&15&18\\\hline
 f_i&6&18&30&24&15&7
\end{array}
\]
The theoretical probability is $P(6\le X<15)=0.75$.
\begin{enumerate}[label=\textbf{(\alph*)}]
\item Find the empirical probability of the stated interval from the given sample. \score{0.5}
\item Calculate the absolute and percentage errors relative to the given theoretical probability. \score{0.5}
\end{enumerate}

\SubsectionBox{6. Download times, grouped data \score{C6-\sml{2}}}
A sample of download times is summarized by the following frequencies.
All times are in seconds; classes are left-inclusive and right-exclusive.
\[
\begin{array}{c|rrrr}
\text{Class}&[0,5)&[5,10)&[10,15)&[15,20)\\\hline
f_i&4&8&6&2
\end{array}
\]
The model gives $P(5\le X<15)=0.80$.
\begin{enumerate}[label=\textbf{(\alph*)}]
\item Find the empirical probability of the stated interval from the given sample. \score{1}
\item Calculate the absolute and percentage errors relative to the given theoretical probability. \score{1}
\end{enumerate}

\SubsectionBox{7. Ticket queues, histogram data \score{C6-\sml{2}}}
The histogram is the only record of a sample of queue times. Bar heights
are frequencies; times are in minutes, with left-inclusive, right-exclusive classes.
\begin{center}
\HistogramGraph{(0,3),(3,6),(6,9),(9,12)}{6,12,9,3}
{0.90\linewidth}{4cm}{blue!25}{3}{3}
\end{center}\vspace{-4mm}
The theoretical probability is $P(3\le X<9)=0.75$.
\begin{enumerate}[label=\textbf{(\alph*)}]
\item Find the empirical probability of the stated interval from the given sample. \score{1}
\item Calculate the absolute and percentage errors relative to the given theoretical probability. \score{1}
\end{enumerate}

Use the sample-variance denominator $n-1$ throughout these four problems.
For each statistic, calculate absolute error $E=|\text{sample value}-\text{theoretical value}|$
and percentage error $100E/|\text{theoretical value}|$. Use unrounded values in all error calculations.
For grouped data and histograms, use class midpoints to estimate the statistics.

\SubsectionBox{8. Sensor readings, raw data \score{C7-\sml{1}}}
Five sensor readings in degrees Celsius are $9,11,12,13,15$.
The theoretical model has $\mu=12$ degrees Celsius and
$\sigma^2=4$ square degrees Celsius.
\begin{enumerate}[label=\textbf{(\alph*)}]
\item Calculate the sample mean, sample variance, and sample standard deviation from the given data. \score{0.5}
\item Calculate the absolute and percentage errors of the mean, variance, and standard deviation relative to the theoretical values. \score{0.5}
\end{enumerate}

\SubsectionBox{9. Items sold, frequency data \score{C7-\sml{1}}}
The numbers of items sold in six short sales periods are summarized below.
The theoretical model has $\mu=5$ items and $\sigma^2=5$ items$^2$.
\[
\begin{array}{c|rrr}x_i&2&5&8\\\hline f_i&1&3&2\end{array}
\]
\begin{enumerate}[label=\textbf{(\alph*)}]
\item Calculate the sample mean, sample variance, and sample standard deviation from the given data. \score{0.5}
\item Calculate the absolute and percentage errors of the mean, variance, and standard deviation relative to the theoretical values. \score{0.5}
\end{enumerate}

\SubsectionBox{10. Cycling times, grouped data \score{C7-\sml{2}}}
Cycling times in minutes are grouped as follows, using left-inclusive,
right-exclusive classes. The theoretical mean is 18 minutes and variance
is 20 minutes$^2$.
\[
\begin{array}{c|rrr}
\text{Time (min)}&[10,14)&[14,18)&[18,22)\\\hline
f_i&1&2&3
\end{array}
\]
\begin{enumerate}[label=\textbf{(\alph*)}]
\item Calculate the sample mean, sample variance, and sample standard deviation from the given data. \score{1}
\item Calculate the absolute and percentage errors of the mean, variance, and standard deviation relative to the theoretical values. \score{1}
\end{enumerate}

\SubsectionBox{11. Crate masses, histogram data \score{C7-\sml{2}}}
A histogram gives a sample of crate masses in kilograms. Bar heights are
frequencies; classes are left-inclusive and right-exclusive. The model
has $\mu=5$ kg and $\sigma^2=2$ kg$^2$.
\begin{center}
\HistogramGraph{(2,4),(4,6),(6,8)}{2,3,1}
{0.90\linewidth}{4.5cm}{orange!30}{2}{1}
\end{center}
\begin{enumerate}[label=\textbf{(\alph*)}]
\item Calculate the sample mean, sample variance, and sample standard deviation from the given data. \score{1}
\item Calculate the absolute and percentage errors of the mean, variance, and standard deviation relative to the theoretical values. \score{1}
\end{enumerate}

\end{multicols}
\end{document}
